\section{Numerical Experiments}
\label{sec:experiments}

The experiments follow the comparison hierarchy dictated by the method's
scope.  First, we compare directly with EFGP-CG, the reference line whose
Fourier model our method accelerates; EFGP already reports broad comparisons
with related large-scale KRR approximations \cite{greengard2025efgp}.  We
report the inverse and active-box EigenPro branches separately and count
full-grid EigenPro as the limiting member of the latter.  Second, we study
box/rank and data-set variation under this same two-family organization, then
repeat the comparison for squared exponential and Mat\'ern kernels.  Third,
we briefly place the best member of each proposed family beside data-space
Nystr\"om and RPCholesky, EFGP+Jacobi, and general-purpose EigenPro baselines.
Finally, we remove model-construction differences by giving every solver
exactly the same hashed \(A\betavec=\bvec\).

The spectral methods in these comparisons form one proposed family rather
than mutually unrelated algorithms.  The active inverse and active-eigenpair
routes spend a controlled budget on a score-selected Fourier box.  The
full-grid EigenPro route is the limiting member obtained when localization is
removed and the spectral correction is constructed on the complete Fourier
grid.  Its stronger result is therefore evidence for the spectral
preconditioning principle; the active default is the lower-setup,
scalability-oriented member of the same family.

The end-to-end timing unit is each method's own algorithmic setup plus solve.
It excludes common data I/O, backend creation, host-to-device staging, and
prediction, which is measured separately.  We report all successful
time--accuracy outcomes rather than suppressing an algorithm because its RMSE
differs from another method by a fixed percentage.  Fixed-system solver totals
include active-set selection when applicable, preconditioner construction,
and CG/PCG; their one shared Fourier construction is reported separately.
Throughout this section, \emph{EFGP-CG} means unpreconditioned CG in the
reference EFGP implementation of Greengard et al.\
\cite{greengard2025efgp}.

\subsection{Setup}
\label{subsec:exp-setup}

The synthetic task uses \(x_i\sim\mathrm{Unif}([0,1]^2)\) and
\[
f_\star(x)=\sin(2\pi x_1)+\tfrac12\cos(2\pi x_2)
+\tfrac15\sin\!\bigl(2\pi(x_1+x_2)\bigr).
\]
Training targets are \(y_i=f_\star(x_i)+\varepsilon_i\), where
\(\varepsilon_i\sim\mathcal N(0,0.3^2)\).  The four exact per-size artifacts
use the same generation procedure and seeds (20260421 for training and 1 for
testing), so their training arrays are consistent nested prefixes.  Each file
contains a noise-free test array of size \(N/4\); formal reporting evaluates
the first 2.5 million test rows.  The retained original configuration sweep
uses the same target function, noise level, and seed convention.  Its timings
are explicitly labelled as the original regime map and are never pooled with
the later five-repeat replay or fixed-system medians.

The real-data task uses ground-class LiDAR points from the 2018 Winnebago USGS
project \cite{usgs3dep_lpc,usgs_winnebago_2018_laz}. Horizontal coordinates
are shifted and rescaled by a common factor and elevations are standardized
using training-set statistics.  Exact nested train artifacts are used at all
four sizes, and formal evaluation uses the first 2.5 million rows of the fixed
held-out test artifact.

The formal comparison uses the Mat\'ern-\(3/2\) kernel with lengthscale
\(\ell=0.1\):
\[
k_{\mathrm{Mat\acute{e}rn}\text{-}3/2}(x,z)
=\left(1+\frac{\sqrt{3}r}{\ell}\right)
\exp\!\left(-\frac{\sqrt{3}r}{\ell}\right),
\]
where \(r=\|x-z\|_2\).  At Fourier-approximation tolerance \(10^{-5}\), the
normalized support has \(M=35721\) modes
\cite{greengard2025efgp,barnett2024uniform}.  We use unit kernel variance,
\(\lambda=0.1\), and relative residual tolerance \(10^{-7}\).  The
prospective broad-pipeline replay and fixed-system study run on an NVIDIA
A100-SXM4-40GB GPU with one warm-up and five measured repetitions.  The
original family regime map is a retained configuration sweep and is labelled
separately rather than pooled with those medians.  Further settings and the
older squared exponential setup are listed in
\cref{tab:common-experiment-parameters}.
\subsection{Primary comparison with EFGP-CG by proposed family}
\label{subsec:end-to-end-krr}

Our primary comparison does not compress the proposed method into a single
automatic ``default'' row.  For each data set and $N$, we report three
quantities: EFGP-CG, the fastest tested active-inverse configuration, and the
fastest tested active-box EigenPro configuration.  The last family includes
the full-grid limit $B=J_m$.  Thus a full-grid winner is counted as a
success of the proposed spectral family, not as an external baseline.

\Cref{tab:primary-two-family-scale} gives the retained A100 Mat\'ern sweep.
The table shows time and test RMSE together; accuracy is not used as a binary
filter for display.  The two daggered inverse rows are the fastest recorded
inverse attempts but were not selected by the original operation-level
fit/convergence screen.  We retain them to make the time--accuracy tradeoff
visible and do not use them for a headline speed claim.  The exact historical
screen, the resulting regime map, and every original candidate row are
reported in \cref{app:archived-regime-map}.

\begin{table}[H]
\centering
\scriptsize
\setlength{\tabcolsep}{2.5pt}
\caption{Primary A100 comparison with EFGP-CG, separated into the inverse and
active-box EigenPro families.  Times are complete recorded setup-plus-solve
totals in seconds.  Each proposed-family entry gives configuration, total,
and test RMSE.  $J_m$ denotes the full Fourier grid.  A dagger marks a raw
inverse attempt rejected by the original screen; it is shown but not used for
a speed claim.}
\label{tab:primary-two-family-scale}
\resizebox{\textwidth}{!}{%
\begin{tabular}{llrrlrrlrr}
\toprule
& & \multicolumn{2}{c}{EFGP-CG} &
\multicolumn{3}{c}{Inverse family} &
\multicolumn{3}{c}{Active-box EigenPro family} \\
\cmidrule(lr){3-4}\cmidrule(lr){5-7}\cmidrule(lr){8-10}
Data & $N$ & Total & RMSE & Configuration & Total & RMSE & Configuration & Total & RMSE \\
\midrule
\multirow{4}{*}{Synthetic}
& $10^7$ & 4.553 & 0.008249 & $|B|=2601$ & \textbf{0.553} & 0.008280 & $|B|=5329,q=256$ & 1.255 & 0.008281 \\
& $3{\times}10^7$ & 6.742 & 0.005238 & $|B|=2601$ & \textbf{0.792} & 0.005262 & $|B|=10609,q=320$ & 1.600 & 0.005262 \\
& $10^8$ & 10.539 & 0.002954 & $|B|=5329$ & \textbf{1.757} & 0.002968 & $J_m,q=320$ & 2.044 & 0.002969 \\
& $3{\times}10^8$ & 15.328 & 0.001729 & $|B|=5329$ & \textbf{2.862} & 0.001738 & $J_m,q=320$ & 3.345 & 0.001739 \\
\addlinespace[1pt]
\multirow{4}{*}{Winnebago}
& $10^7$ & 5.309 & 0.115140 & $|B|=625$ & 4.550 & 0.115147 & $J_m,q=128$ & \textbf{0.979} & 0.115057 \\
& $3{\times}10^7$ & 8.643 & 0.113434 & $|B|=625$ & 4.684 & 0.113583 & $J_m,q=128$ & \textbf{1.378} & 0.113378 \\
& $10^8$ & 15.462 & 0.111901 & $|B|=5329^{\dagger}$ & 6.984 & 0.150677 & $J_m,q=320$ & \textbf{1.603} & 0.111862 \\
& $3{\times}10^8$ & 27.748 & 0.110611 & $|B|=5329^{\dagger}$ & 8.038 & 0.262925 & $J_m,q=384$ & \textbf{2.901} & 0.110583 \\
\bottomrule
\end{tabular}%
}
\end{table}

The two branches are complementary.  On Synthetic, the inverse family is
fastest at all four displayed sizes and gives $5.36\times$--$8.52\times$
complete-pipeline speedup over EFGP-CG.  On Winnebago, the active-box
EigenPro family is fastest and improves from $5.42\times$ at $10^7$ to
$9.56\times$ at $3\times10^8$.  This is a data-dependent regime change,
not an unexplained discontinuity in one fixed algorithm.  It is precisely why
we report the two families separately instead of allowing an automatic route
switch to appear as one continuous ``ours'' curve.

\iffalse
\subsection{Primary end-to-end comparison with EFGP-CG}
\label{subsec:end-to-end-krr}

We first compare the proposed binned EFGP default directly with EFGP-CG on
Synthetic and Winnebago at four sizes from \(10^7\) to \(3\times10^8\).
Both methods complete all eight cases.  The default's paired-median total-time
speedup ranges from \(3.24\times\) to \(10.75\times\).  The gain is not
confined to one timing component: paired setup speedups range from
\(2.12\times\) to \(3.67\times\), and paired solve speedups range from
\(3.28\times\) to \(17.92\times\).

\begin{table}[H]
\centering
\scriptsize
\setlength{\tabcolsep}{2.4pt}
\caption{Primary A100 comparison with EFGP-CG.  Times are in seconds;
``total'' is the median method-owned algorithmic setup-plus-solve time, while
its two components come from the repeat attaining that median total.  \(S\)
is the median of repeat-wise EFGP-CG/ours total-time ratios, so it need not
equal the ratio of the displayed totals.  Prediction is separate.}
\label{tab:prospective-efgp-scale}
\resizebox{\textwidth}{!}{%
\begin{tabular}{llrrrrrrrrr}
\toprule
& & \multicolumn{3}{c}{EFGP-CG} & \multicolumn{3}{c}{Ours} & &
\multicolumn{2}{c}{Test RMSE} \\
\cmidrule(lr){3-5}\cmidrule(lr){6-8}\cmidrule(lr){10-11}
Data & \(N\) & Setup & Solve & Total & Setup & Solve & Total & \(S\) & CG & Ours \\
\midrule
\multirow{4}{*}{Synthetic}
& \(10^7\)              & 0.157 & 4.447 & 4.604 & 0.070 & 1.041 & 1.111 & 4.138 & 0.008249 & 0.008280 \\
& \(3{\times}10^7\)    & 0.489 & 6.382 & 6.871 & 0.157 & 1.456 & 1.613 & 4.317 & 0.005239 & 0.005262 \\
& \(10^8\)              & 1.496 & 9.504 & 11.000 & 0.464 & 1.459 & 1.923 & 5.694 & 0.002954 & 0.002969 \\
& \(3{\times}10^8\)    & 4.483 & 11.688 & 16.171 & 1.346 & 2.366 & 3.712 & 4.356 & 0.001730 & 0.001740 \\
\addlinespace[1pt]
\multirow{4}{*}{Winnebago}
& \(10^7\)              & 0.166 & 5.119 & 5.285 & 0.072 & 1.553 & 1.625 & 3.238 & 0.115140 & 0.115057 \\
& \(3{\times}10^7\)    & 0.465 & 8.412 & 8.877 & 0.163 & 2.316 & 2.479 & 3.548 & 0.113406 & 0.113351 \\
& \(10^8\)              & 1.532 & 14.368 & 15.900 & 0.447 & 1.118 & 1.565 & 10.167 & 0.111989 & 0.111949 \\
& \(3{\times}10^8\)    & 4.949 & 23.840 & 28.788 & 1.349 & 1.334 & 2.682 & 10.754 & 0.090216 & 0.090239 \\
\bottomrule
\end{tabular}%
}
\end{table}

The advantage remains pronounced for \(N\geq10^8\), where paired total
speedup is \(4.36\times\)--\(10.75\times\).  At
\(N=3\times10^8\), Synthetic setup and solve fall from \(4.483\) and
\(11.688\) seconds to \(1.346\) and \(2.366\) seconds, respectively, for a
\(4.36\times\) total speedup.  On Winnebago the same comparison is \(28.788\)
versus \(2.682\) seconds, a \(10.75\times\) speedup.  All eight proposed rows
remain within one percent of the corresponding full-grid spectral reference;
relative RMSE changes against that reference range from \(-0.07\%\) to
\(+0.57\%\).  Thus the large-sample acceleration retains comparably high
predictive quality while making the small remaining accuracy differences
explicit.

\fi

\subsection{Family-specific parameter and data-set variation}
\label{subsec:formal-robustness}

We next keep the two branches separate while changing their internal
parameters and the data set.  Candidate configurations are compared by total
time, not by iteration count.  Test RMSE is carried from the same row and is
reported alongside time.  This distinction matters because a larger inverse
box or a higher EigenPro rank can shorten the solve while increasing the
one-time construction cost.

\begin{table}[H]
\centering
\footnotesize
\setlength{\tabcolsep}{4pt}
\caption{Within-family parameter variation for Synthetic Mat\'ern at
$N=10^7$.  Totals include the recorded binned setup, preconditioner
construction, and PCG.  The lowest iteration count is not necessarily the
lowest total.}
\label{tab:two-family-parameter-variation}
\begin{tabular}{llrrr}
\toprule
Family & Configuration & Iter. & Total (s) & Test RMSE \\
\midrule
\multirow{5}{*}{Inverse}
& $|B|=625$  & 594 & 0.976 & 0.0082805 \\
& $|B|=961$  & 408 & 0.711 & 0.0082805 \\
& $|B|=1369$ & 325 & 0.617 & 0.0082805 \\
& $|B|=2601$ & 174 & \textbf{0.553} & 0.0082805 \\
& $|B|=5329$ &  88 & 1.127 & 0.0082805 \\
\addlinespace[1pt]
\multirow{4}{*}{Active-box EigenPro}
& $|B|=2601,q=192$  & 510 & 1.274 & 0.0082804 \\
& $|B|=5329,q=256$  & 416 & \textbf{1.255} & 0.0082805 \\
& $|B|=10609,q=320$ & 365 & 1.289 & 0.0082804 \\
& $J_m,q=256$       & 389 & 1.527 & 0.0082806 \\
\bottomrule
\end{tabular}
\end{table}

The inverse optimum is the intermediate box $|B|=2601$: enlarging the box to
5329 halves the iteration count again but doubles construction sufficiently
to make the total slower.  The active-box EigenPro totals are flatter across
the tested localized boxes, with $|B|=5329,q=256$ narrowly best.  Hence the
default-setting argument is family-specific; a single top-$k$/rank pair cannot
stand in for both operations.

Data geometry can also reverse the family ordering.  At $N=3\times10^7$ with
the same Mat\'ern kernel, Synthetic favors the inverse branch
(0.792 versus 1.600 seconds), whereas Winnebago favors the full-grid member of
the active-box EigenPro branch (1.378 versus 4.684 seconds).  The Fourier
diagonal and score alone do not determine this switch: the sample-dependent
Toeplitz term changes with the input distribution.  The revised one-click
artifact therefore writes a separate inverse row and a separate
active-box/full-grid minimum for every parameter and data-set case; it never
lets an automatic route switch masquerade as stability of one method.

\subsection{Kernel comparison under the same two-family definition}
\label{subsec:kernel-family-comparison}

The previous mechanism table mixed score fractions, eigenvalues, and a
particular preconditioner without making the compared totals clear.  We
replace it with a direct kernel experiment.  For each kernel and data set,
\cref{tab:kernel-two-family} uses the family-optimal configuration selected by
the procedure of \cref{subsec:end-to-end-krr}: one inverse result and the
minimum over localized and full-grid EigenPro results.  Every numeric column
has the same meaning as in the primary table.

\begin{table}[H]
\centering
\footnotesize
\setlength{\tabcolsep}{3.5pt}
\caption{Two-family comparison at $N=3\times10^7$.  Total is complete
recorded setup plus solve in seconds; speedup is EFGP-CG total divided by the
displayed family total.  Full-grid EigenPro is denoted by $J_m$.}
\label{tab:kernel-two-family}
\resizebox{\textwidth}{!}{%
\begin{tabular}{lllrrrr}
\toprule
Data & Kernel & Family/configuration & Total & Test RMSE & Speedup & EFGP-CG \\
\midrule
\multirow{4}{*}{Synthetic}
& SE & Inverse, $B=J_m$ & \textbf{0.211} & 0.001071 & 14.70 & 3.098 \\
& SE & Active-box EigenPro, $J_m,q=128$ & 0.810 & 0.001073 & 3.82 & 3.098 \\
& Mat\'ern & Inverse, $|B|=2601$ & \textbf{0.792} & 0.005262 & 8.52 & 6.742 \\
& Mat\'ern & Active-box EigenPro, $|B|=10609,q=320$ & 1.600 & 0.005262 & 4.21 & 6.742 \\
\addlinespace[1pt]
\multirow{4}{*}{Winnebago}
& SE & Inverse, $B=J_m$ & \textbf{0.212} & 0.374291 & 13.55 & 2.868 \\
& SE & Active-box EigenPro, $J_m,q=64$ & 0.499 & 0.374291 & 5.74 & 2.868 \\
& Mat\'ern & Inverse, $|B|=625$ & 4.684 & 0.113583 & 1.85 & 8.643 \\
& Mat\'ern & Active-box EigenPro, $J_m,q=128$ & \textbf{1.378} & 0.113378 & 6.27 & 8.643 \\
\bottomrule
\end{tabular}%
}
\end{table}

For SE, the grid has only 1225 modes and full inversion wins on both data
sets.  For Mat\'ern, the grid has 35721 modes: localization makes the inverse
branch best on Synthetic, while the global eigenpair correction is much
better on Winnebago.  The conclusion is therefore not that one
preconditioner dominates universally, but that the proposed construction
contains two useful computational regimes whose best member depends on
kernel spectrum and data geometry.

\iffalse
\subsection{Parameter sensitivity and default-setting stability}
\label{subsec:formal-robustness}

The scale experiment prospectively selects the Synthetic
\(N=3\times10^7\) Mat\'ern case as the controlled target: its EFGP-CG solve
requires 5968 iterations, inside the predeclared 3000--6000 target window.
Starting from this frozen case, we vary one factor at a time.  The same six
complete KRR pipelines are rerun at every point; method parameters are not
retuned using the new timings.  For lengthscale and data-set changes, a frozen
score prefix may only be shortened deterministically when its enclosing box
would exceed the unchanged budget.

\begin{table}[H]
\centering
\footnotesize
\setlength{\tabcolsep}{4.2pt}
\caption{One-at-a-time robustness at the frozen \(N=3\times10^7\) target.
Times are method-owned setup plus solve in seconds.  The last column is the
paired-median EFGP-CG/default speedup.  All proposed rows satisfy the declared
broad RMSE and \(R^2\) usability range.}
\label{tab:formal-robustness}
\begin{tabular}{llrrrr}
\toprule
Axis & Value & EFGP-CG & Default & Test RMSE & Speedup \\
\midrule
\multirow{3}{*}{\(\lambda\)}
& 0.01 & 17.044 & 2.632 & 0.005603 & 6.48 \\
& 0.1  &  6.919 & 1.613 & 0.005262 & 4.29 \\
& 1.0  &  2.962 & 1.160 & 0.004196 & 2.55 \\
\addlinespace[1pt]
\multirow{3}{*}{\(\ell\)}
& 0.05 & 6.853 & 2.349 & 0.009366 & 2.92 \\
& 0.1  & 6.919 & 1.613 & 0.005262 & 4.29 \\
& 0.2  & 4.655 & 1.195 & 0.002859 & 3.90 \\
\addlinespace[1pt]
\multirow{3}{*}{Box budget}
& 4096  & 6.847 & 0.997 & 0.005262 & 6.87 \\
& 8192  & 6.755 & 1.598 & 0.005262 & 4.23 \\
& 16384 & 6.924 & 1.711 & 0.005262 & 4.05 \\
\addlinespace[1pt]
\multirow{2}{*}{Data}
& Synthetic & 7.034 & 1.599 & 0.005262 & 4.40 \\
& Winnebago & 8.860 & 2.341 & 0.113351 & 3.78 \\
\bottomrule
\end{tabular}
\end{table}

The default remains faster than EFGP-CG at every declared value, by
\(2.55\times\)--\(6.87\times\).  This supports robustness of the acceleration
relative to the EFGP-CG line; it does not assert that the localized default is
the fastest member of the proposed family at every point.  In particular, a
full-grid spectral correction can be preferable when the difficult
eigenspace is not sufficiently localized.  That outcome is a success of the
family's extreme full-grid member, not evidence against spectral
preconditioning.

The archived box--rank scan supplies the configuration study used before the
formal timing.  On the Winnebago Mat\'ern problem at \(N=3\times10^8\), smaller
boxes improve steadily as rank grows but remain slower than the full-grid
route.  On the full grid, increasing \(q\) from 192 to 448 reduces iterations
from 794 to 272, whereas total time is minimized at the intermediate
\(q=384\).  Thus neither the largest rank nor the smallest iteration count is
an appropriate default by itself; the selected configuration balances box
coverage, construction, and solve.  The complete 15-row values are retained
in \cref{tab:box-rank-sweep}.

\begin{figure}[H]
    \centering
    \includegraphics[width=0.72\linewidth]{picture/group_c_boxeig_parameter_sweep.png}
    \caption{Archived full-grid/active-box rank sweep used to study default
    stability.  The formal campaign transfers only the selected box and rank;
    it reruns all timing and does not inherit archived timing values.}
    \label{fig:archived-box-rank-sweep}
\end{figure}

\subsection{Mechanism check}
\label{subsec:mechanism-check}

To interpret the total-time trends, we check the two observations used by the method:
the score should identify a limited part of the Fourier grid, and EigenPro
should reduce the large eigenvalues inside a chosen box. These checks use saved
diagnostics at \(N=10^8\) and \(\lambda=0.1\); they are not additional timing
runs.

For each kernel, order the indices by decreasing \(\rho_j\) and measure the
fraction of \(\sum_j\rho_j\) in the leading indices. We also examine the USGS
Mat\'ern block with \(|B|=2601\) and \(q=256\). The first two rows of
\cref{tab:mechanism-check} compare score concentration; the last two compare
the measured block spectrum before and after the EigenPro operation.

\begin{table}[H]
\centering
\small
\caption{Two checks at \(N=10^8\). Score rows give the fraction of total
\(\rho\); spectrum rows use the USGS Mat\'ern block with \(|B|=2601\) and
\(q=256\). The transformed values are calculated from the measured spectrum.}
\label{tab:mechanism-check}
\begin{tabular}{lrr}
\toprule
Leading indices & SE score fraction & Mat\'ern score fraction \\
\midrule
256 & \(0.998\) & \(0.869\) \\
512 & \({>}0.99999\) & \(0.942\) \\
\addlinespace[2pt]
USGS Mat\'ern box & Raw \(A_{BB}\) & After \(P_{\mathrm{eig}}(B,256)\) \\
\cmidrule(lr){1-3}
Largest eigenvalue & \(2.16\times10^7\) & \(1.00\) \\
Condition number & \(2.16\times10^8\) & \(4.75\times10^1\) \\
\bottomrule
\end{tabular}
\end{table}

The score is much more concentrated for the squared exponential kernel. On
the tested Mat\'ern box, the EigenPro operation reduces the block condition
number from \(2.16\times10^8\) to \(47.5\). This checks the operation on
\(A_{BB}\), but not the approximation on \(T\) or the coupling \(A_{BT}\);
the complete solves below test those remaining effects.

\subsection{Additional archived EFGP-CG regime map}
\label{subsec:exp-scalability}

The original six-size sweep is retained because it extends the primary
EFGP-CG comparison down to \(N=10^6\) and includes both squared exponential
and Mat\'ern kernels.  \Cref{tab:se-scalability,tab:matern-scalability} report
the EFGP-CG baseline and a representative configuration saved by that sweep.
These archived ratios are not combined with the formal A100 ratios above.  A
full-grid configuration is written as
\(P_{\mathrm{inv}}(J_m)\) or \(P_{\mathrm{eig}}(J_m,q)\); it is not treated
as a separate method. Speedup is the EFGP-CG total time divided by the
candidate total time. It is therefore a complete-pipeline speedup, including
the setup route used by each archived pipeline.  Both time and RMSE are shown
so that the speed--accuracy tradeoff remains visible.

\begin{table}[H]
\centering
\footnotesize
\setlength{\tabcolsep}{3.5pt}
\caption{Archived squared exponential EFGP scalability results (\(M=1225\)).
Each pair compares EFGP-CG with the representative saved candidate.
\(T_{\mathrm{tot}}\) is total time in seconds, RMSE uses the same test set,
and speedup compares the recorded complete pipelines.}
\label{tab:se-scalability}
\begin{tabular}{lllrrrr}
\toprule
Data & \(N\) & Method / configuration & Iter. &
\(T_{\mathrm{tot}}\) & Test RMSE & Speedup \\
\midrule
\multirow{12}{*}{Synthetic}
& \multirow{2}{*}{\(10^6\)} & EFGP-CG & 1115 & 1.23 & \(4.90{\times}10^{-3}\) & 1.00 \\
& & \(P_{\mathrm{inv}}(B),\ |B|=625\) & 10 & 0.05 & \(4.91{\times}10^{-3}\) & 25.35 \\
& \multirow{2}{*}{\(3{\times}10^6\)} & EFGP-CG & 1495 & 1.64 & \(2.91{\times}10^{-3}\) & 1.00 \\
& & \(P_{\mathrm{inv}}(B),\ |B|=625\) & 16 & 0.07 & \(2.91{\times}10^{-3}\) & 23.34 \\
& \multirow{2}{*}{\(10^7\)} & EFGP-CG & 2051 & 2.35 & \(1.64{\times}10^{-3}\) & 1.00 \\
& & \(P_{\mathrm{inv}}(J_m)\) & 1 & 0.11 & \(1.70{\times}10^{-3}\) & 21.99 \\
& \multirow{2}{*}{\(3{\times}10^7\)} & EFGP-CG & 2613 & 3.10 & \(9.85{\times}10^{-4}\) & 1.00 \\
& & \(P_{\mathrm{inv}}(J_m)\) & 1 & 0.21 & \(1.07{\times}10^{-3}\) & 14.70 \\
& \multirow{2}{*}{\(10^8\)} & EFGP-CG & 3344 & 4.74 & \(5.72{\times}10^{-4}\) & 1.00 \\
& & \(P_{\mathrm{inv}}(J_m)\) & 1 & 0.56 & \(6.88{\times}10^{-4}\) & 8.46 \\
& \multirow{2}{*}{\(3{\times}10^8\)} & EFGP-CG & 2777 & 9.80 & \(3.42{\times}10^{-4}\) & 1.00 \\
& & \(P_{\mathrm{inv}}(J_m)\) & 1 & 1.55 & \(5.65{\times}10^{-4}\) & 6.31 \\
\addlinespace[1pt]
\multirow{12}{*}{USGS}
& \multirow{2}{*}{\(10^6\)} & EFGP-CG & 717 & 0.79 & \(3.95{\times}10^{-1}\) & 1.00 \\
& & \(P_{\mathrm{inv}}(J_m)\) & 1 & 0.06 & \(3.95{\times}10^{-1}\) & 12.28 \\
& \multirow{2}{*}{\(3{\times}10^6\)} & EFGP-CG & 1030 & 1.15 & \(3.90{\times}10^{-1}\) & 1.00 \\
& & \(P_{\mathrm{inv}}(J_m)\) & 1 & 0.08 & \(3.90{\times}10^{-1}\) & 14.07 \\
& \multirow{2}{*}{\(10^7\)} & EFGP-CG & 1495 & 1.70 & \(3.82{\times}10^{-1}\) & 1.00 \\
& & \(P_{\mathrm{inv}}(J_m)\) & 1 & 0.11 & \(3.82{\times}10^{-1}\) & 15.93 \\
& \multirow{2}{*}{\(3{\times}10^7\)} & EFGP-CG & 2371 & 2.87 & \(3.74{\times}10^{-1}\) & 1.00 \\
& & \(P_{\mathrm{inv}}(J_m)\) & 1 & 0.21 & \(3.74{\times}10^{-1}\) & 13.55 \\
& \multirow{2}{*}{\(10^8\)} & EFGP-CG & 3541 & 4.83 & \(3.65{\times}10^{-1}\) & 1.00 \\
& & \(P_{\mathrm{inv}}(J_m)\) & 1 & 0.56 & \(3.65{\times}10^{-1}\) & 8.66 \\
& \multirow{2}{*}{\(3{\times}10^8\)} & EFGP-CG & 5611 & 10.75 & \(3.56{\times}10^{-1}\) & 1.00 \\
& & \(P_{\mathrm{inv}}(J_m)\) & 1 & 1.81 & \(3.56{\times}10^{-1}\) & 5.95 \\
\bottomrule
\end{tabular}
\end{table}

For the squared exponential kernel, \(M=1225\) is small enough that the
full-grid inverse is usually fastest. It reduces PCG to one iteration, but it
cannot by itself reduce the \(O(N)\) Fourier setup. Consequently, the speedup
decreases at the largest \(N\), where setup forms a larger fraction of total
time. The two smaller synthetic rows prefer \(|B|=625\), showing that a partial
inverse can be cheaper when CG is not yet the dominant cost.

On the squared exponential synthetic task, the proposed test RMSE is visibly
higher at the two largest sizes, although its absolute value remains below
\(6\times10^{-4}\).  We retain these rows and report both RMSE columns rather
than hiding that tradeoff behind a binary accuracy decision.

\begin{table}[H]
\centering
\footnotesize
\setlength{\tabcolsep}{3.5pt}
\caption{Archived Mat\'ern-\(3/2\) EFGP scalability results (\(M=35721\)). Each pair
compares EFGP-CG with the representative saved candidate. Times are
seconds, and speedup compares the recorded complete pipelines.}
\label{tab:matern-scalability}
\begin{tabular}{lllrrrr}
\toprule
Data & \(N\) & Method / configuration & Iter. &
\(T_{\mathrm{tot}}\) & Test RMSE & Speedup \\
\midrule
\multirow{12}{*}{Synthetic}
& \multirow{2}{*}{\(10^6\)} & EFGP-CG & 1806 & 1.94 & \(1.87{\times}10^{-2}\) & 1.00 \\
& & \(P_{\mathrm{inv}}(B),\ |B|=1369\) & 130 & 0.27 & \(1.89{\times}10^{-2}\) & 7.23 \\
& \multirow{2}{*}{\(3{\times}10^6\)} & EFGP-CG & 2791 & 6.95 & \(1.31{\times}10^{-2}\) & 1.00 \\
& & \(P_{\mathrm{inv}}(B),\ |B|=1369\) & 188 & 0.39 & \(1.31{\times}10^{-2}\) & 17.87 \\
& \multirow{2}{*}{\(10^7\)} & EFGP-CG & 4163 & 4.55 & \(8.25{\times}10^{-3}\) & 1.00 \\
& & \(P_{\mathrm{inv}}(B),\ |B|=2601\) & 174 & 0.55 & \(8.28{\times}10^{-3}\) & 8.24 \\
& \multirow{2}{*}{\(3{\times}10^7\)} & EFGP-CG & 5968 & 6.74 & \(5.24{\times}10^{-3}\) & 1.00 \\
& & \(P_{\mathrm{inv}}(B),\ |B|=2601\) & 262 & 0.79 & \(5.26{\times}10^{-3}\) & 8.52 \\
& \multirow{2}{*}{\(10^8\)} & EFGP-CG & 8831 & 10.54 & \(2.95{\times}10^{-3}\) & 1.00 \\
& & \(P_{\mathrm{inv}}(B),\ |B|=5329\) & 187 & 1.76 & \(2.97{\times}10^{-3}\) & 6.00 \\
& \multirow{2}{*}{\(3{\times}10^8\)} & EFGP-CG & 11494 & 15.33 & \(1.73{\times}10^{-3}\) & 1.00 \\
& & \(P_{\mathrm{inv}}(B),\ |B|=5329\) & 246 & 2.86 & \(1.74{\times}10^{-3}\) & 5.36 \\
\addlinespace[1pt]
\multirow{12}{*}{USGS}
& \multirow{2}{*}{\(10^6\)} & EFGP-CG & 1804 & 1.95 & \(1.24{\times}10^{-1}\) & 1.00 \\
& & \(P_{\mathrm{eig}}(J_m,128)\) & 131 & 0.63 & \(1.23{\times}10^{-1}\) & 3.09 \\
& \multirow{2}{*}{\(3{\times}10^6\)} & EFGP-CG & 2876 & 3.11 & \(1.18{\times}10^{-1}\) & 1.00 \\
& & \(P_{\mathrm{eig}}(J_m,128)\) & 203 & 0.75 & \(1.18{\times}10^{-1}\) & 4.13 \\
& \multirow{2}{*}{\(10^7\)} & EFGP-CG & 4914 & 5.31 & \(1.15{\times}10^{-1}\) & 1.00 \\
& & \(P_{\mathrm{eig}}(J_m,128)\) & 341 & 0.98 & \(1.15{\times}10^{-1}\) & 5.42 \\
& \multirow{2}{*}{\(3{\times}10^7\)} & EFGP-CG & 7796 & 8.64 & \(1.13{\times}10^{-1}\) & 1.00 \\
& & \(P_{\mathrm{eig}}(J_m,128)\) & 534 & 1.38 & \(1.13{\times}10^{-1}\) & 6.27 \\
& \multirow{2}{*}{\(10^8\)} & EFGP-CG & 13539 & 15.46 & \(1.12{\times}10^{-1}\) & 1.00 \\
& & \(P_{\mathrm{eig}}(J_m,320)\) & 264 & 1.60 & \(1.12{\times}10^{-1}\) & 9.65 \\
& \multirow{2}{*}{\(3{\times}10^8\)} & EFGP-CG & 22466 & 27.75 & \(1.11{\times}10^{-1}\) & 1.00 \\
& & \(P_{\mathrm{eig}}(J_m,384)\) & 320 & 2.90 & \(1.11{\times}10^{-1}\) & 9.56 \\
\bottomrule
\end{tabular}
\end{table}

The Mat\'ern grid is almost thirty times larger. The archived rows select
moderate inverse boxes for the synthetic task and the full-grid spectral
family member for the largest USGS cases. The explanation cannot come from
\(D\) or \(\rho_j\) alone: for fixed
kernel parameters, grid, \(N\), and \(\lambda\), these quantities agree, while
the sample-dependent Toeplitz matrix \(G\) changes with the input points.

The Mat\'ern baseline and selected test RMSE values differ by about one percent
or less in every row.

Legacy timing components and the full box--rank sweep are reported in
\cref{tab:timing-breakdown,tab:box-rank-sweep}.

\fi

\subsection{Brief comparison with other complete KRR pipelines}
\label{subsec:other-krr}

We next report the other complete KRR pipelines as a secondary comparison.
The table shows the largest scale, where scalability is most consequential;
the complete \(10^7\)--\(3\times10^8\) matrix is retained in the experiment
artifact.  No fixed relative-accuracy cutoff is used to remove a successful
row: raw training total and test RMSE are reported together.  The full-grid
spectral row is the proposed family's high-budget extreme, whereas Nystr\"om
and RPCholesky are data-space KRR algorithms and Jacobi is an EFGP baseline.
EFGP-CG remains the primary reference and is already reported in
\cref{tab:primary-two-family-scale}.

\begin{table}[H]
\centering
\footnotesize
\setlength{\tabcolsep}{4.2pt}
\caption{Secondary end-to-end KRR comparison at \(N=3\times10^8\).
Training total is method-owned setup plus solve in seconds.  ``Resource
limit'' is the prospectively declared memory outcome, not a timed failure.}
\label{tab:other-krr-300m}
\begin{tabular}{lrrrr}
\toprule
& \multicolumn{2}{c}{Synthetic} & \multicolumn{2}{c}{Winnebago} \\
\cmidrule(lr){2-3}\cmidrule(lr){4-5}
Method & Train total & Test RMSE & Train total & Test RMSE \\
\midrule
Ours (binned default) & \textbf{3.712} & 0.001740 & \textbf{2.682} & 0.090239 \\
Full-grid spectral family member & 6.442 & 0.001730 & 6.253 & 0.090214 \\
EFGP+Jacobi & 5.270 & 0.001730 & 31.406 & 0.090216 \\
Nystr\"om KRR & 14.797 & 0.023004 & 14.793 & 0.275328 \\
RPCholesky KRR & \multicolumn{2}{c}{Resource limit} & \multicolumn{2}{c}{Resource limit} \\
\bottomrule
\end{tabular}
\end{table}

At \(3\times10^8\), the binned default is the fastest successful route on both
data sets.  The full-grid family member is also successful: it preserves the
EFGP-level RMSE while substantially improving on EFGP-CG, and is
\(1.74\times\) and \(2.34\times\) slower than the cheaper binned default on
Synthetic and Winnebago, respectively.  Nystr\"om exposes a different
speed--accuracy tradeoff at smaller sizes; at the largest size shown here it
is both slower and less accurate.  RPCholesky completes the \(10^7\) cases,
but its declared memory requirement exceeds the available budget from
\(3\times10^7\) onward, so we make no large-\(N\) timing claim for it.

\paragraph{General-purpose EigenPro line.}
EigenPro2 and EigenPro3 attack the data-space kernel-training problem with
GPU-oriented stochastic updates and Nystr\"om spectral information
\cite{ma2019eigenpro,abedsoltan2023toward}.  They therefore do not share the
EFGP Fourier system or its residual-based stopping rule.  We keep their
comparison separate from the \(10^7\)--\(3\times10^8\) complete-pipeline
matrix and use time to an EFGP-CG validation target.  The retained legacy
EigenPro3 experiment uses \(9\times10^4\) training, \(10^4\) validation, and
\(10^3\) test points.  For Mat\'ern-\(3/2\), the \(p=30000\) model does not
reach the target within 30 epochs, whereas the proposed Fourier family reaches
it in 0.50 seconds.  For the Gaussian kernel, EigenPro3 reaches the target in
158.24 seconds versus 0.29 seconds for the proposed route.  EigenPro2 was also
considered as the predecessor in this line, but no matched, reproducible
time-to-target row is retained in the formal artifact; consequently we do not
turn its missing row into a numerical speedup or a convergence claim.

\subsection{Controlled fixed-system comparison}
\label{subsec:matched-confirmation}

To isolate the solver contribution, we select the Synthetic
\(N=3\times10^7\) Mat\'ern case and construct the Fourier weights, Toeplitz
generator, regularization arrays, and right-hand side exactly once.  All six
methods receive the same zero initial iterate and immutable system, whose
SHA-256 identifier begins \texttt{12c0a9811168}.  The common \(A,\bvec\)
construction takes \(0.606\) seconds and is excluded from every solver total.
After one warm-up, five shuffled measured repetitions synchronize each timing
boundary.  For active methods the total includes fresh active-set selection,
preconditioner construction, and solve; for CG it is the solve alone.

\begin{table}[H]
\centering
\footnotesize
\setlength{\tabcolsep}{3.4pt}
\caption{Fixed-system A100 comparison at \(N=3\times10^7\).  Components and
totals are median-total representatives in seconds; \(S_{\mathrm{CG}}\) is
the median of matched-repeat CG/method total ratios.  All five measured
repetitions converge with independently recomputed relative residual below
\(10^{-7}\).}
\label{tab:fixed-system-formal}
\begin{tabular}{lrrrrrr}
\toprule
Method & Iter. & Select & Build & Solve & Total & \(S_{\mathrm{CG}}\) \\
\midrule
CG & 6185 & 0 & 0 & 6.525 & 6.525 & 1.000 \\
Jacobi & 1385 & 0 & 0.0005 & 1.761 & 1.761 & 3.715 \\
Ours (frozen default) & 495 & 0.0048 & 0.678 & 0.798 & 1.481 & 4.421 \\
Active EigenPro & 495 & 0.0048 & 0.714 & 0.809 & 1.527 & 4.272 \\
Full-grid spectral family member & 479 & 0 & 0.519 & 0.757 & \textbf{1.276} & \textbf{5.128} \\
Active inverse & 54 & 0.0049 & 3.734 & 0.146 & 3.884 & 1.667 \\
\bottomrule
\end{tabular}%
\end{table}

The frozen default reduces the total from \(6.525\) to \(1.481\) seconds, a
\(4.42\times\) paired-median speedup over CG, and wins all five matched
repetitions.  Because \(A\) and \(\bvec\) are identical, this confirms that
the gain over the EFGP-CG line is not only a faster setup effect.  More
importantly, the full-grid family member reaches \(1.276\) seconds and a
\(5.13\times\) speedup.  We interpret this as the high-budget limiting success
of the proposed spectral family: removing localization spends less selection
logic and captures the global difficult subspace more completely on this
particular system.  The default remains the deployment-oriented route because
its localized construction is designed to scale when a full-grid spectral
setup becomes less attractive.  The active inverse illustrates the opposite
tradeoff: its 54-iteration, 0.146-second solve is the shortest, but the
3.734-second factor construction raises its total to 3.884 seconds.

\FloatBarrier
